\documentclass{article}
\usepackage{amsmath}
\begin{document}
\title{The addressed cave observation has a radio path, but needs a phase budget}
\maketitle

\paragraph{Source and physical boundary.}
The earlier coupled cave study (\#781) gives the installer a
blocked-move outcome and needs the courier to act on that result
at its next movement cut. The branch positions used here are
installer at 13 or 42 and courier at 16. Using the same sampled
SubT medium and chosen 72 m hard range, the selected source
broker delivers a named one-byte report from the installer to the
courier from either position, if dispatch runs and the named
capacity permits. The old report from tile 18 to the base
receiver at 39 is separately received in every case.
This does not turn authored safe movement outcomes into measured
geometry. The framing uses a one-byte message to represent a
binary choice, not one byte of physical airtime or authenticated
sender proof. The civil fixture has exactly the named known
participants and their authorized topology.

\paragraph{Reception limit is indexed by the sending configuration.}
The selected \texttt{DispatchMessages} passes the transmitting
member's radio configuration to the channel function. That
function applies its capacity both to the sender's transmit
ledger and to the destination's current receive ledger.
A preceding one-byte monitor record at the courier can therefore
block the installer report under the installer's cap of one,
even when the courier's own configuration is larger. Conversely
in this source it can be accepted when the courier's nominal
transmit configuration is smaller but the installer's cap
admits it. The holds are only those actually incurred, not
declared packet lengths alone. Exactly one second old counts
against the tested received budget. Eighty of 320 selected
cases disagree with a rule that indexes the receive
limit by the courier's own configuration. This is a source
interpretation for the fixed parameter arms, not a recommendation
to change senders or a global service claim about every packet.
Where all radios share the same cap the two expressions coincide.
Separately, a registered member whose pose read fails can leave
the bit queued through the first native dispatch, even with
both actual communicating endpoints ready. The later recovery
in that arm is provided by the model.

\paragraph{Exact ideal phase cutoff.}
The intended movement resolution occurs at integer time \(m\)
within a normalized 1024-tick slot. Qualification and enqueue
take \(q\); decode takes \(\rho\). The receiver's action cut is
1024. If enqueue is proved to precede dispatch at equality,
the first eligible tick is \(u=m+q\). Otherwise strict order
makes it \(u=m+q+1\). With one periodic integer dispatch grid of
period \(P\), if all phases are possible and all due attempts
can otherwise run, the latest first dispatch is \(u+P-1\).
Thus every phase fits exactly when
\[
 u+P-1+\rho\le1024.
\]
This excludes roster stalls, new interference during the wait,
and clock misalignment. It is a proof about the declared
integer schedule, not a measured timer period. In particular a
move that resolves at the decision boundary gives no positive
slack for forwarding even if both radio routes are feasible.
Checking the same reachability in a new position does not supply
the missing slack. The numerical enumeration compares a direct
phase scan, and the selected native replay distinguishes receipt
and capacity from schedule alignment. No independent cave
mission or full simulator timeline has been run.
\end{document}
